%% Methods subsection for the reframed manuscript.
%% Two hero analyses underpinning the dispersion-controlled topological readout:
%%   (a) covariance-matched Gaussian dispersion null  (scripts/17_dispersion_null_and_stability.py)
%%   (b) bootstrap directional robustness             (scripts/17_dispersion_null_and_stability.py)
%% with the circular-coordinate negative control      (scripts/16_lineage_circular_coords.py)
%% All quantitative statements trace to the locked evidence in
%% results/foundation/ and manuscript/audit/RESCUE_STRATEGY_2026-06-10.md.
%% No traversal/cycling claim is made anywhere in this file.

\subsection*{Covariance-matched Gaussian dispersion null}

The central methodological control of this study isolates genuine
non-Gaussian topological structure from the trivial dependence of
$H_1$ persistence on the overall spread of the point cloud. A more
dispersed cluster of points admits larger empty regions and therefore
larger $H_1$ bars even in the complete absence of any organised
loop-like geometry; a rise in max $H_1$ is therefore not, by itself,
evidence of topological reorganisation. To separate the two, we
compare each observed value against the max $H_1$ of a multivariate
Gaussian that is matched to the data in mean and covariance --- and
hence in dispersion --- but that contains no structure beyond a single
unimodal ellipsoidal blob.

For a given cohort group we take the cell$\times$PC matrix on the
top-30 PCA coordinates ($n_{\text{PCs}} = 30$, consistent with the
rest of the analysis) and estimate its empirical mean $\mu$ and
covariance $\Sigma$. We then draw synthetic point clouds of identical
size from $\mathcal{N}(\mu, \Sigma)$ by Cholesky factorisation,
$\Sigma + 10^{-9} I = L L^{\top}$, generating each null cloud as
$\mu + Z L^{\top}$ with $Z$ a standard-normal matrix of the same shape
as the subsample. For every synthetic cloud we recompute the
Vietoris--Rips $H_1$ persistence diagram with \texttt{ripser} 0.6
\citep{bauer2021,tralie2018} under the same Euclidean metric and
$\mathbb{F}_2$ coefficients used for the real data, and record the
maximum finite-bar persistence $\max(\text{death} - \text{birth})$.
This yields a null distribution of max $H_1$ that carries exactly the
data's first- and second-order moments but, by construction, none of
its higher-order structure. We summarise the comparison as the ratio
$\text{ratio} = \widetilde{H_1^{\text{real}}} \,/\,
\widetilde{H_1^{\text{gauss}}}$ of the median real to median null
persistence. A $\text{ratio} \approx 1$ indicates that max $H_1$ is
acting as a pure dispersion proxy at that condition, whereas
$\text{ratio} > 1$ (real exceeding the dispersion-matched null)
indicates genuine non-Gaussian topological structure beyond spread.

This control is decisive at the level of \emph{when} structure
appears, not merely whether persistence is large. Under osimertinib in
the PC9 cell line, the dispersion-matched ratio is $0.95$ at baseline
(D0) and $0.91$ at D3 --- at or below unity, i.e.\ max $H_1$ is
indistinguishable from a dispersion artefact early in treatment ---
and then rises to $1.31$ at D7 and $1.42$ at D14 as drug-tolerant
persisters emerge, demonstrating structure in excess of dispersion
specifically under sustained drug exposure. The same contrast holds in
the EGFR-mutant PDX model YU-006, where the ratio is $1.21$ in the
untreated condition and $1.95$ in the osimertinib-residual condition.
In the erlotinib PC9 cohort the ratios remain at or below unity at
every stage ($0.93$, $0.99$ and $0.70$ at D0, D9 and D11), so this
cohort shows no excess structure and we treat it as a specificity
control rather than a positive finding (see Discussion). In the
patient cohorts (Maynard et al.) the ratio exceeds unity already at the
treatment-naive stage ($1.52$, with $1.69$ and $1.51$ at the
persister and progressive-disease stages); because the excess is
present in naive tumours it reflects baseline tumour heterogeneity
rather than drug induction, and we accordingly scope the
drug-induction claim to the controlled osimertinib cell-line and PDX
systems.

\subsection*{Bootstrap directional robustness}

Point estimates of max $H_1$ are sensitive to which cells are sampled
(coefficient of variation $7$--$19\%$ across the osimertinib and PDX
groups), so we do not rest any claim on a single numerical value.
Instead we test whether the \emph{ordering} of conditions --- drug
above baseline, and real above the dispersion-matched null --- is
stable under resampling. For each cohort group we draw $30$
independent subsamples of $1{,}200$ cells without replacement (or the
full group size where smaller) under a fixed random seed
(\texttt{numpy} default generator, seed $42$) and compute max $H_1$
for each. We report the bootstrap median, the $90\%$ interval as the
$5$th--$95$th percentiles, and the coefficient of variation; for the
Gaussian comparison, dispersion-matched null clouds are drawn from the
subsamples (with an additional richer draw on a representative
subsample) and pooled into the null distribution against which the
real values are compared.

Two robustness quantities support the central claims. First, the
fraction of (real, null) pairs in which the real max $H_1$ exceeds the
dispersion-matched null measures how reliably structure exceeds
spread: this is $31\%$ and $18\%$ at osimertinib D0 and D3 (real
not exceeding null), rising to $100\%$ at D7 and $98\%$ at D14; it is
$86\%$ for YU-006 untreated and $99\%$ for YU-006 residual. Second,
directional robustness across condition pairs measures how reliably
the drug condition exceeds baseline: osimertinib D7 exceeds D0 and
osimertinib D14 exceeds D0 in $100\%$ of bootstrap subsample pairs,
and the YU-006 residual condition exceeds untreated in $97\%$ of
pairs. The reported real max $H_1$ medians and $90\%$ intervals are
$1.41\,[1.31, 1.60]$ (osi D0), $1.82\,[1.82, 2.15]$ (osi D3),
$2.67\,[2.45, 3.04]$ (osi D7) and $3.78\,[2.99, 3.78]$ (osi D14) for
the cell line, and $2.79\,[2.23, 4.03]$ (untreated) and
$5.46\,[3.58, 7.34]$ (residual) for YU-006. Because the claims are
phrased as orderings holding in $97$--$100\%$ of resamples rather than
as differences between point estimates, they are robust to the
disclosed subsample instability.

\subsection*{Circular-coordinate negative control}

To guard against over-interpreting the persistent $H_1$ loop as a
trajectory that cells move around, we ran an explicit construct-validity
check whose purpose is to \emph{bound} the claim, not to extend it. For
each osimertinib timepoint we computed circular coordinates
$\theta \in [0, 1)$ for every cell from the most-persistent $H_1$
cohomology class using the de Silva--Morozov--Vejdemo-Johansson
algorithm (greedy-permutation landmarks with harmonic smoothing of the
lifted $\mathbb{Z}/p$ cocycle, odd prime field $p = 47$, up to $400$
landmarks) as implemented in \texttt{dreimac}
\citep{perea2023}. If cells were distributed around the loop, these
coordinates would spread across the circle; instead they collapse to a
single angular concentration, with a Rayleigh resultant length
$R \geq 0.965$ across all cohorts, far above the $R = 0.881$ obtained
for a synthetic control loop occupied at $12\%$ of its circumference.
Using Watermelon lineage barcodes (GSE150949), the angular dispersion
(circular variance, $1 - \bar{R}$) of clones with at least three
co-sampled cells was compared to size-matched random cell groups over
$1{,}000$ permutations; clonally related cells were \emph{tighter} than
random, not more spread. These results establish that the structure is
a sparse persistent loop occupied by a minority of cells; we therefore
make no claim that cells traverse it, and present this analysis as a
negative control demonstrating the absence of such traversal.

%% Reproducibility: analyses (a) and (b) are implemented in
%% scripts/17_dispersion_null_and_stability.py and the negative control in
%% scripts/16_lineage_circular_coords.py; summaries are written to
%% results/foundation/dispersion_null_summary.json and
%% results/foundation/circular_coord_summary.json. Seed 42 throughout.
